\section{Selection Matrix}

The matrix maps task type to algorithm family, CLI surface, evidence
requirement, and audit artifact. Construction methods provide baselines, search
methods explore alternatives, exact methods can provide bounded certificates,
multi-objective methods expose trade-offs, and audit methods verify declared
artifact contracts.

\begin{table}[htbp]
\centering
\small
\begin{tabular}{p{0.20\linewidth}p{0.22\linewidth}p{0.24\linewidth}p{0.22\linewidth}}
\toprule
Need & Preferred family & Required evidence & Audit handoff \\
\midrule
Fast baseline & Construction & L0/L1 fixture and smoke tests & RPLAN with constructor metadata \\
Improve candidate & Local/LNS search & Validity-preserving move trace & RPLAN/RCTX run summary \\
Bounded claim & Exact optimization & Bound, witness, or solve report & Certificate and manifest \\
Distributional context & Ensemble/SMC & Protocol and sampled-plan record & Exported plans and context \\
Trade-off choice & Pareto selection & Frontier and selection rule & Selected-plan package \\
Artifact verification & Audit verifier & Declared profile pass/fail & Certificate bundle \\
\bottomrule
\end{tabular}
\caption{Method-family selection by claim need and audit handoff.}
\end{table}

The matrix is conservative by design. If the claim requires a bound, a search
heuristic can provide useful candidates but should not be described as exact.
If the claim requires distributional context, a single optimized plan is an
example, not a sample. If the claim is legal or policy-facing, the workflow
must carry enough evidence to support interpretation without letting the audit
certificate stand in for that interpretation.

\subsection*{Preflight Questions}

Before selecting a workflow, the operator should answer four questions:
What claim is the output meant to support? What evidence level is required for
that claim? Which CLI surface or crate currently produces the relevant plan or
certificate artifact? Which RPLAN/RCTX sidecars must be retained so another
reader can inspect the run?
